\section{Execution: build readiness, distributed placement, and repository-local windows}
\label{sec:execution}

\Cref{sec:protocol} states what must happen before what. This section states
what that costs in wall time, where the tasks may run, and which history each
window ranges over when the packages live in several repositories. Its first
subsection is informative in the sense of \spec{13.11}: the propositions bound
durations of schedules that respect \cref{def:readiness}, and no conforming
implementation is required to attain them. Its second and third subsections
transcribe the two optional profiles of \spec{28} and \spec{27} as definitions
and state the one invariance each profile promises. Neither extends the
theorems of \cref{sec:guarantees}, which are stated for a single history.

\subsection{Build readiness durations}
\label{sec:exec-durations}

\spec{19.2a} promises a happens-before relation and no duration; \spec{13.11}
nevertheless says what each relation buys when nothing else limits the run.
Throughout, package $i$ has a build of duration $b_i$ and a publication of
duration $p_i$, both fixed; a task starts as soon as its predecessors in the
task graph of \cref{def:taskgraph} have finished and a slot is free; and one
relation is declared on every edge.

\begin{proposition}[the longest path, \spec{13.11}]
\label{prop:exec-longest}
With unlimited build and publication slots a run lasts exactly the longest
path $L$ of its task graph, measured in task durations. A weaker relation
never lengthens it.
\end{proposition}

\begin{proof}
The task graph is acyclic (\cref{prop:readiness-acyclic}), so its tasks can be
taken in a topological order. With unlimited slots a task starts at the
greatest finish time of its predecessors, or at $0$ if it has none, and by
induction along that order its finish time is the longest path ending at it.
The run ends at the greatest finish time, which is the longest path of the
graph. A weaker relation contributes a subset of the edges
(\cref{prop:readiness-acyclic}), so every path of its task graph is a path of
the stronger one's, and the longest is no longer.
\end{proof}

\begin{proposition}[a chain, \spec{13.11}]
\label{prop:exec-chain}
Let packages $1, \dots, n$ form a chain in which package $k+1$ consumes package
$k$, under one relation throughout. With unlimited slots the run lasts
\[
  \sum_{i=1}^{n} (b_i + p_i) \quad\text{under \code{publish}},\qquad
  \max_{1 \le k \le n}\Bigl(\sum_{i=1}^{k} b_i + \sum_{i=k}^{n} p_i\Bigr)
  \quad\text{under \code{build}},
\]
\[
  \max_{1 \le k \le n}\Bigl(b_k + \sum_{i=k}^{n} p_i\Bigr)
  \quad\text{under \code{none}}.
\]
\end{proposition}

\begin{proof}
Under \code{publish}, $\exbuild{k+1}$ waits for $\expub{k}$, so the $2n$ tasks
form one path and the sum is its length. Under \code{build}, the builds form
the path $\exbuild{1} \to \dots \to \exbuild{n}$ and $\exbuild{k}$ ends at
$B_k = \sum_{i \le k} b_i$; the publication of $k$ waits for $\exbuild{k}$ and
for $\expub{k-1}$, so its finish $E_k$ satisfies $E_k = \max(B_k, E_{k-1}) +
p_k$ with $E_0 = 0$. By induction $E_k = \max_{j \le k}\bigl(B_j +
\sum_{i=j}^{k} p_i\bigr)$: the case $k = 1$ is $B_1 + p_1$, and the step is
$\max\bigl(B_k + p_k,\ \max_{j < k}(B_j + \sum_{i=j}^{k-1} p_i) + p_k\bigr)$.
Under \code{none} every build starts at $0$ and $\exbuild{k}$ ends at $b_k$;
the same recurrence with $b_k$ in place of $B_k$ gives the third form. In each
case the run lasts $E_n$ by \cref{prop:exec-longest}.
\end{proof}

\begin{corollary}[equal durations]
\label{cor:exec-equal}
With $b_i = b$ and $p_i = p$ for all $i$, the three forms are $n(b+p)$,
$\max(nb + p,\ b + np)$ and $b + np$. \code{build} saves $(n-1)\min(b,p)$
over \code{publish}, and \code{none} saves a further $(n-1)\max(b-p, 0)$, so
the second saving exists only where builds outlast publications.
\end{corollary}

\begin{proof}
Under \code{build} the maximand $kb + (n-k+1)p$ is affine in $k$, so the
maximum is at $k = 1$ or $k = n$. Under \code{none} the maximand $b + (n-k+1)p$
is greatest at $k = 1$. The differences follow by cases on $b \ge p$.
\end{proof}

\begin{proposition}[fan-out and fan-in, \spec{13.11}]
\label{prop:exec-fan}
For one provider $0$ and consumers $i \ge 1$, each consuming $0$ and nothing
else, the run lasts
\[
  b_0 + p_0 + \max_i (b_i + p_i),\qquad
  b_0 + \max_i\bigl(\max(b_i, p_0) + p_i\bigr),\qquad
  \max_i\bigl(\max(b_i,\ b_0 + p_0) + p_i\bigr)
\]
under \code{publish}, \code{build} and \code{none} respectively. For providers
$i \ge 1$ of one consumer $c$, each consumed by $c$ and nothing else, it lasts
\[
  \max_i (b_i + p_i) + b_c + p_c,\qquad
  \max\bigl(\max_i b_i + b_c,\ \max_i (b_i + p_i)\bigr) + p_c,\qquad
  \max\bigl(b_c,\ \max_i (b_i + p_i)\bigr) + p_c.
\]
\end{proposition}

\begin{proof}
Fan-out. Under \code{publish}, consumer $i$'s build starts at $b_0 + p_0$ and
its publication follows its build, giving $b_0 + p_0 + b_i + p_i$; take the
maximum. Under \code{build}, its build starts at $b_0$ and ends at $b_0 + b_i$,
and its publication waits for that and for $\expub{0}$ at $b_0 + p_0$, giving
$b_0 + \max(b_i, p_0) + p_i$. Under \code{none}, its build ends at $b_i$ and its
publication waits for that and for $b_0 + p_0$. Fan-in. Under \code{publish},
$c$'s build starts when the last provider has published, at $\max_i (b_i +
p_i)$. Under \code{build}, it starts at $\max_i b_i$ and ends $b_c$ later, and
$c$'s publication waits for that and for every provider's publication. Under
\code{none}, $c$'s build ends at $b_c$ and only its publication waits for the
providers. Each is the longest path of \cref{prop:exec-longest}.
\end{proof}

\begin{proposition}[one publication lane, \spec{13.11}]
\label{prop:exec-lane}
Where one owner's publications are serialised (\spec{27.2}, \spec{28.6}) in
the order $1, \dots, n$ of \cref{def:order}, let $e_k$ be the instant the build
of package $k$ ends. Under every relation the run lasts
\[
  \max_{1 \le k \le n}\Bigl(e_k + \sum_{i=k}^{n} p_i\Bigr),
\]
which is at least $\sum_{i} p_i$ whatever the relation, and a weaker relation
only lowers the $e_k$.
\end{proposition}

\begin{proof}
Publication $k$ starts at $\max(e_k, E_{k-1})$, since it waits for its own
build and for the lane, and the lane order agrees with \cref{def:order}, so
every provider of $k$ has published by $E_{k-1}$. The recurrence is that of
\cref{prop:exec-chain} with $e_k$ in place of $B_k$, and the same induction
gives the maximum. The term $k = 1$ is $e_1 + \sum_i p_i \ge \sum_i p_i$. For
the last claim, add the lane edges $\expub{k-1} \to \expub{k}$ to the task
graph; they follow the publish order, which is a topological order of the
workspace graph, so the key of \cref{prop:readiness-acyclic} still increases
along every edge and the graph stays acyclic. With a fixed lane order every
task starts at the greatest finish of its predecessors, and by induction along
a topological order each finish time under a subset of the edges is at most
the finish time under the full set; $e_k$ is such a finish time.
\end{proof}

\begin{proposition}[budgets, \spec{13.11}]
\label{prop:exec-budget}
Let there be $m_b$ build slots and $m_p$ publication slots, and write $W_b =
\sum_i b_i$, $W_p = \sum_i p_i$ and $L$ for the longest path of the task
graph. No schedule finishes before $\max(L,\ W_b/m_b,\ W_p/m_p)$. A schedule
that never leaves a slot of a stage idle while a task of that stage is ready
finishes within $L + W_b/m_b + W_p/m_p$, which is within three times the best
possible. Under budgets a weaker relation can lengthen such a schedule, but
never the best one.
\end{proposition}

\begin{proof}
The lower bound: a schedule contains every path, so it lasts at least $L$; it
performs $W_b$ units of build work on $m_b$ slots, so at least $W_b/m_b$; and
likewise $W_p/m_p$. The upper bound: let the schedule end at $T$ and let $t_1$
be a task finishing at $T$. For $j \ge 1$, while $t_j$ has predecessors, let
$t_{j+1}$ be one that finishes last among them; this gives a chain $t_q \to
\dots \to t_1$ whose total duration is at most $L$. Consider an instant of
$[0, T)$ at which no task of the chain is running. It lies between the finish
of some $t_{j+1}$ and the start of $t_j$ (or before the start of $t_q$), and
at that instant $t_j$ is ready, since its latest predecessor has finished, and
has not started; by hypothesis every slot of $t_j$'s stage is then busy. The
instants at which all $m_b$ build slots are busy total at most $W_b/m_b$, and
those at which all $m_p$ publication slots are busy at most $W_p/m_p$. Hence $T
\le L + W_b/m_b + W_p/m_p$, and each summand is at most the lower bound. For
the last claim, a schedule feasible under a stronger relation is feasible under
a weaker one, so the best weaker schedule is no longer; that a particular
greedy schedule can be longer is \cref{ex:exec-anomaly}.
\end{proof}

The bound is Graham's argument for list scheduling under precedence
constraints~\cite{graham66,graham69}, applied to two stages with separate
budgets; \spec{13.11} does not claim it tight for this shape of graph.

\begin{example}[Graham's anomaly, \spec{13.11}]
\label{ex:exec-anomaly}
Take two build slots and one publication slot, packages $A = (2, 3)$ and $B =
(1, 2)$ where $B$ consumes $A$, and an unrelated $C = (8, 3)$, with ready tasks
taken in the order $A, B, C$. Under \code{build}, $B$'s build is not ready at
$0$: $A$ builds in $[0,2)$ and $C$ in $[0,8)$ on the two slots; $B$ builds in
$[2,3)$; $A$ publishes in $[2,5)$, $B$ in $[5,7)$ once $A$ has published, and
$C$ in $[8,11)$ once its build ends. The run lasts $11$. Under \code{none},
$B$'s build is ready at $0$ and takes the second slot: $A$ builds in $[0,2)$
and $B$ in $[0,1)$, so $C$ builds only in $[1,9)$; $A$ publishes in $[2,5)$,
$B$ in $[5,7)$ and $C$ in $[9,12)$. The run lasts $12$: the consumer's early
build took the slot the long build would have had. A list that preferred $C$
to $B$ at $0$ would finish at $11$ under \code{none} as well, which is the
point: the length depends on the list, not on the relation, and it is why
\spec{19.2a} promises a happens-before relation and no duration.
\end{example}

\begin{example}[an infrastructure chain]
\label{ex:exec-chain}
Let the infrastructure leg of \cref{ex:terraform} be consumed by a backend
consumed in turn by a frontend, with $(b, p)$ in minutes of $(4, 6)$ for the
infrastructure, $(9, 3)$ for the backend and $(12, 2)$ for the frontend. By
\cref{prop:exec-chain}, under \code{publish} the run lasts $10 + 12 + 14 = 36$
minutes; under \code{build} it lasts $\max(4 + 11,\ 13 + 5,\ 25 + 2) = 27$;
under \code{none} it lasts $\max(4 + 11,\ 9 + 5,\ 12 + 2) = 15$. The
publications still run in order in every case, and the frontend's publication
still ends the run.
\end{example}

\subsection{Distributed execution}
\label{sec:exec-distributed}

\spec{28} is an optional execution profile that changes no message syntax and
asserts no speedup. The definitions below transcribe its obligations; the one
proposition is the invariance it promises.

\begin{definition}[nodes and roles, \spec{28.1}]
\label{def:exec-roles}
An execution node is a machine or runner executing the same release engine and
configuration schema; it is distinct from a repository peer. Every node
supports two roles, and one engine provides both. The default role is
\emph{orchestrator}. An explicit \emph{worker} refuses release initiation,
including indirect initiation from a delegated hook, and executes authorised
tasks in an existing run; it never substitutes a locally recomputed plan,
acquires competing release ownership, dispatches to another pool, or
independently finalises a release. An orchestrator-capable node that accepts
a delegated task has only worker authority for it. The node on which CI starts
the release is the run's orchestrator, of which there is one per run; it owns
lock acquisition, the fixed planning input, task allocation, result admission
and release finalisation, and may execute tasks locally under the same rules
as a worker. Workers may be added to every history mode; execution links are
not fleet links and alter neither source discovery nor release scope.
\end{definition}

\begin{definition}[lock, snapshot, plan, assign, \spec{28.3}]
\label{def:exec-input}
Before creating a write-capable run the orchestrator identifies every
participating repository, acquires all release locks in one stable total
order, verifies ownership, and revalidates discovery under those locks; a
changed participant set releases the partial acquisition and acquires the
complete set again. Discovery and compatibility inspection before locking are
read-only, and an unsafe lock bypass is refused. The run input $\exI$ contains
the fixed repository-qualified heads and complete relevant history, the release
records of the authoritative stores as compared under the locks
(\cref{def:records-lock}), the graph, the resolved semantic configuration, the
explicit release options, and the withdrawal inventory and receipts of
\spec{26}. The semantic release plan is the pure function $\exPlan(\exI)$.
Node placement configuration is not semantic input: worker count, placement,
completion order, wall-clock time, run identifiers and transport branch names
do not change $\exPlan(\exI)$. A plan digest is computed over a canonical
serialisation of semantic input and plan content, and transient execution
identity is bound separately. The orchestrator derives a task graph from the
plan, including preparation, tests, builds, transfers, publications and
recording gates; a task is placed only on a compatible node and only when its
exact prerequisites are available; the schedule need not be deterministic.
Where a package's stages may be placed is execution policy declared beside
\cref{def:readiness} and, like it, no part of the plan. Every reconciliation
that writes a shared manifest or lockfile is an explicit task under the
orchestrator's authority, serialised per shared file; each admitted result is
one prepared input state, builds fan out only from an admitted state, and
every build binds the state it consumed. An assignment identifies the run,
plan digest, task, attempt, ownership generation, repository-qualified source
object identifiers, effective configuration and command environment,
toolchain and platform requirements, input receipts, permitted writes and
outputs; workers authenticate the assigning orchestrator and the assignment's
integrity; duplicate delivery is recognised; only one attempt is authorised
to perform a given external side effect; and a task receipt is execution
evidence, never a release record.
\end{definition}

\begin{proposition}[placement invariance]
\label{prop:exec-placement}
For a fixed input $\exI$, the plan $\plan{\st}$, the versions
$\vers{\st}{\cdot}$ and the publish order of \cref{def:order} do not depend on
the placement of tasks, on the number of nodes, or on the order in which tasks
complete.
\end{proposition}

\begin{proof}
By \cref{thm:determinism} the plan, the versions and the planned publish
order are functions of $(\st, \Head, \mathrm{config})$ alone. $\exI$ fixes all
three: \st{} is the set of release records compared under the locks
(\cref{def:records-lock}), \Head{} is the fixed repository-qualified head, and
the configuration is the run's resolved semantic configuration. Nothing a node
contributes reaches those arguments: the node-startup settings of \spec{28.2}
(role, capacity, worker links) do not override package policy, planned
versions, commands or inputs; the build readiness relation and the placement
policy are execution policy outside the plan (\cref{def:readiness},
\cref{def:exec-input}); and placement, node count and completion order are
excluded from $\exI$ by \cref{def:exec-input}. A function of fixed arguments
takes one value. Vectors 3 and 5 of \spec{28.8} are this statement.
\end{proof}

\begin{definition}[Git synchronisation branches, \spec{28.4}]
\label{def:exec-branches}
The orchestrator provisions isolated task worktrees and temporary coordination
branches named \code{dispat-worker-<id>-<workinfo>}, created atomically under
\code{refs/heads/}, where \code{<id>} is the node the branch is addressed to
and \code{<workinfo>} carries a date, a kind label and at least 128 bits of
randomness. The name is an untrusted routing hint that lets a node list the
branches addressed to it; the authenticated manifest, not the name,
identifies node, run, task and attempt, and a node rejects a manifest whose
binding disagrees with where it was found. Each acknowledged checkpoint
identifies an exact full object identifier, and a consumer reads only from
that object, never from a tip that can still move; a tip fetched but not read
or verified has not been seen. Ref updates and deletions check the expected
previous object identifier and attempt ownership, so that an update racing
another party's update fails instead of overwriting it. Transport credentials
do not authorise a worker to update native release branches, tags or fleet
settlement refs. A source snapshot may travel without the history behind it,
as a parentless commit or bare tree, because a task's commands may not read
history from a transport checkout. Transport commits are never release intent,
a window boundary, a successful release or an admitted native head; they are
neither merged into a release branch nor made ancestors of release tags, and
only validated release-file changes enter native recording through the
ordinary transaction. Artefact provenance binds the exact source revision and
admitted release-file state used for the build.
\end{definition}

\begin{definition}[outputs as first-class inputs, \spec{28.5}]
\label{def:exec-outputs}
A provider's revision does not make its ignored or generated outputs
available on another machine, so a conforming implementation transports the
outputs a local-output dependency needs, declared ignored and generated files
included; it neither reclones and rebuilds the provider closure on every
consumer node nor pins the connected graph to one worker. A successful build
emits a manifest binding its outputs to the run, plan, task, attempt,
repository-qualified source revision, admitted release-file state, package and
version, toolchain and platform, and input and output digests. Identity is by
content: consumers verify the authenticated manifest, content integrity and
compatibility before their stage becomes ready; installation is atomic at the
task-input boundary; a completed build publishes one admitted output set for
its exact attempt; and a missing, stale, incomplete or incompatible output
fails the prerequisite and is never substituted with registry contents or a
new build under a different input. Bytes travel through the coordination
branch or through immutable content-addressed bundles whose references and
digests are in it; a reference alone is not a transfer. For package edge $B
\to A$ under the \code{build} relation, readiness is the task precedence
$\exbuild{A} \to \mathrm{export}(A) \to \mathrm{verify/materialise}(A
\text{ at } B\text{'s node}) \to \exbuild{B}$; under \code{publish} it
requires $A$'s publication and every applicable recording gate, which local
bytes never replace; under \code{none} the edge carries no output and the two
builds are independent tasks. A provider a consumer reaches only across
\code{none} edges is no input of that consumer and is not prepared for it:
the input closure ends at a \code{none} provider. Task inputs include every
required build dependency, including packages needing no new release and
build-only edges outside the propagation kinds; the task graph is acyclic,
and an unschedulable cycle is an execution-configuration error, never a
reason to change release intent. Verified outputs may be reused across runs
only under a new admission receipt bound to the current run, plan and task
attempt.
\end{definition}

\begin{definition}[publication, recording and recovery, \spec{28.6}, \spec{28.2}]
\label{def:exec-publication}
The orchestrator alone admits results and authorises publication, after the
ordinary test, dependency, settlement and relevant-input checks; it serialises
publication and recording within one owner, a single-history monorepo
included, while isolated builds within that owner may overlap and independent
owners keep their publication concurrency. A worker never moves release
branches, creates release tags, settles peer links or finalises the fleet.
Authorisation to publish is explicit, single-use and bound to one task
attempt: it is issued only after the package's \code{beforePublish} hook has
completed on the executing node and after the relevant-input and lock-ownership
checks, the publisher does not start its command before it has observed the
authorisation for its own attempt, and an authorisation is never reused by
another attempt or for another effect. An authorisation may be bounded in
time, with a start instant after which the publisher does not begin and a
deadline at which the executing node ends the command itself; the bound
limits when an effect may begin and how long it runs, is not an expiry of the
release lock, and authorises no takeover. The orchestrator retains and
verifies lock ownership for the lifetime of all authorised effects, and
assignments and results are bound to its ownership generation. After lock
loss no new effect starts: no further assignment and no further
authorisation. Recording an effect authorised before the loss is not a new
effect and is still attempted, as the create-only record of
\cref{def:tagging}, because a publication left unrecorded is a wedge the next
owner can clear only by \cref{def:reconcile}, and a create-only record cannot
overwrite what a successor wrote. A timed-out worker is not proof that its
publisher stopped: before ownership is handed over, cancellation is
acknowledged with the phase the attempt stopped in, or a fence is enforced
where the effect occurs, or the outcome is reconciled; a generation field
alone fences no registry command, and no takeover expires ownership while the
previous process can still act. Capacity returns only on evidence: an
in-flight attempt retains its reservation until completion, acknowledged
cancellation or safe fencing, a timeout alone frees nothing, and an
assignment no node has claimed holds no capacity and may be withdrawn and
offered again. A worker that fails before publication is retried under a new
attempt without discharging the package; a publication that may have
succeeded is reconciled by \cref{def:reconcile} before another attempt, since
neither a missing reply nor a missing tag proves failure; only the durable
source release record discharges the obligation. Locks are released only after
authorised operations have quiesced or been fenced, in reverse acquisition
order; if that cannot be proved the run fails holding its exclusion and no
record of the uncertain package, which the next run treats as pending. An
attempt never authorised to perform an external effect is fenced by revoking
its coordination ref; an authorised publisher is not. An abandoned run is
findable from its lock, and recovery revokes the run's coordination refs,
settles every authorised attempt without a result, and only then removes the
lock.
\end{definition}

\begin{example}[the required example, \spec{28.7}]
\label{ex:exec-required}
In one large JavaScript repository let $\code{assets} \Rightarrow \{\code{ui},
\code{docs}\} \Rightarrow \code{app}$ denote task precedence: \code{ui} and
\code{docs} consume \code{assets}, and \code{app} consumes both; the package
edges of \cref{def:graph} point the other way. Worker $A$ builds \code{assets}
once and exports its required outputs. Workers $B$ and $C$ receive verified
outputs and build \code{ui} and \code{docs} concurrently. The assigned
\code{app} worker receives both output sets and builds \code{app}.
Publication and source recording then obey \cref{def:order} and the owner
lane of \cref{def:exec-publication}. Neither $B$ nor $C$ rebuilds
\code{assets}, and the \code{app} worker rebuilds neither closure. The
example works without converting the monorepo to a polyrepo or publishing
private intermediate packages solely for transport.
\end{example}

\begin{remark}[the performance boundary, \spec{28.7}]
\label{rem:exec-boundary}
The required gain is the distribution and reuse of dependent builds with
parallel ready branches; universal speedup is not a semantic guarantee. Git
provides source and result transport; it does not schedule tasks, transfer
omitted build products, guarantee compatible environments or make effects
exactly-once. A strictly serial chain retains its critical path, and transfer
and setup overhead can outweigh saved compute on small jobs. Before an
implementation claims a performance gain it compares the same pinned fixture
locally and on several nodes, reporting cold and warm setup, per-node
resources, transfer bytes and time, task execution counts, total work, wall
time, repetitions and variability, and artefact equivalence, with the fixture
and environment identified; skipped or failed trials are not speedups, and
neither source-only dispatch nor duplicated prerequisite compilation is a
distributed-build success.
\end{remark}

\begin{proposition}[list scheduling under placement, \spec{13.11}]
\label{prop:exec-graham}
Let $W_k$ be the summed duration of a run's command tasks, $L_k$ the duration
of its longest precedence chain, and $s$ the task slots usable at once, the
lesser of the summed node capacities and the applicable run-wide stage
limits. No schedule finishes before $\max(L_k,\ W_k/s)$. When slots are
interchangeable and transfers cost nothing, a placement that never leaves a
usable slot idle while a compatible task is ready finishes within $W_k/s +
(1 - 1/s)\,L_k$, which is at most $(2 - 1/s)$ times the best possible. On a
chain $L_k = W_k$ and no number of nodes helps.
\end{proposition}

\begin{proof}
The lower bound is that of \cref{prop:exec-budget} with one stage. For the
upper bound take the chain $t_q \to \dots \to t_1$ of that proof. At every
instant either a chain task runs or all $s$ slots are busy, and in the latter
case the work done off the chain is at most $W_k - L'$ where $L' \le L_k$ is
the chain's length; so the off-chain time is at most $(W_k - L')/s$ and the
makespan at most $L' + (W_k - L')/s = W_k/s + (1 - 1/s)L' \le W_k/s + (1 -
1/s)L_k$. Each of $W_k/s$ and $L_k$ is at most the optimum, giving the factor.
This is Graham's bound for list scheduling under precedence
constraints~\cite{graham66,graham69}. On a chain every task is on the
longest path, so $L_k = W_k$ and the lower bound is $W_k$ for every $s$.
\end{proof}

Neither premise holds in general: nodes differ in platform and speed, and an
output transfer delays the consumer it feeds. \spec{13.11} therefore states no
bound once transfers count and puts the measurement of
\cref{rem:exec-boundary} in its place. The inequality still says what to
expect: once $s \ge W_k/L_k$ the lower bound is $L_k$, so the chain and not the
pool is the limit, and more nodes can only close the gap between a greedy
schedule and $L_k$, which is below a factor of two.

\paragraph{Implementation status.}
The design history of \SPEC{} dated 2026-09-22 records that the dispat
release engine implements the profile in every history mode over the Git
mailbox transport of \cref{def:exec-branches}: every message is a signed tree
on a branch, every transition one compare-and-swap ref update, a worker
polls, and nothing connects to a worker. Builds and publications run on
workers, a publication is delegated only by an explicit placement, a space
that logs in to its registry publishes on the orchestrator, the orchestrator
is one more node of the pool chosen last, and a provider the run does not
release is prepared once per run when a consumer reads its outputs. Outputs
travel as result trees on the branches, with no bundle service, which is
vector 17 of \spec{28.8}. Not implemented, so their conditional vectors do
not apply: reuse of verified outputs across runs (vector 20) and the rollback
profile of \spec{26} (vector 26); the durable mailbox ref and the
history-free snapshot that \spec{28.4} permits are likewise unimplemented, so
a snapshot descending from the planned head carries the repository's whole
history to a mailbox that starts empty. Departures from the profile as
written: the platform check of \spec{28.2} runs before dispatch for the
packages a run releases and, for a prepared provider, at its placement; and
the \code{native-recording-or-lock} class is attached where a distributed run
loses its lock and not yet at the older recording and lock sites. The same
history records that the owed windows of \spec{13.3} and \diag{E201}
(\cref{def:owed-commit}) are not yet implemented, so a consumer that sits out
the run in which its provider releases the owed commit is still stranded in
that engine. What was measured, on which fixtures, is reported in
\cref{sec:evaluation}.

\subsection{Repository-local windows}
\label{sec:exec-polyrepo}

\spec{27} plans one package graph whose packages are owned by several linked
Git repositories. It changes where history and release records are read and
nothing else: the grammar, the bump lattice, the windows, holds,
cancellation, corrections, channel rules, the dependency graph and the
partial-publication rules of \cref{sec:protocol} are unchanged. This
subsection defines the objects the profile adds and states what it promises;
it does not extend the single-history theorems of \cref{sec:guarantees} to
several histories.

\begin{definition}[fleet snapshot and commit identity, \spec{27.2}]
\label{def:exec-fleet}
Before parsing history the engine resolves one immutable map from repository
identity to full commit object identifier, the fixed fleet snapshot; every
participating repository is an initialised, non-shallow checkout at its
pinned commit, and a missing, shallow, moved or ambiguous participant is
\diag{E330}. A commit identity is a pair (repository, full object
identifier); two equal identifier byte strings in different repositories are
incomparable. Commit dates, author dates, modification times, tag creation
dates and fetch order create no ancestry or precedence. The engine reads each
repository's reachable DAG once for the snapshot, preserves native parent
order, and manufactures no synthetic fleet DAG or parent edge. Ancestry,
changed paths, direct pending windows, cancellation barriers, corrections and
\code{Reverts} are evaluated in the repository that owns the commit.
\end{definition}

\begin{definition}[repository-local windows, \spec{27.5}, \spec{27.6}, \spec{27.11}]
\label{def:exec-windows}
For a source repository $r$ with fixed head $h_r$, write $\exRch{r}{x}$ for
the ancestors-or-self of a commit $x$ of $r$, and
\[
  \exafter{r}{x} \;=\; \exRch{r}{h_r} \setminus \exRch{r}{x},
  \qquad
  \exafter{r}{\NONE} \;=\; \exRch{r}{h_r}.
\]
For a consumer package $D$ and a repository $r$, let $s_{D,r}$ be the
revision of $r$ that $D$'s last stable release incorporated and $b_{D,r}$ the
revision of $r$ that $D$'s newest release of any channel incorporated, each
$\NONE$ where no such release exists. $D$'s train and fresh windows over $r$
are
\[
  \exWin{D}{r} \;=\; \exafter{r}{s_{D,r}},
  \qquad
  \exWinf{D}{r} \;=\; \exafter{r}{s_{D,r}} \cap \exafter{r}{b_{D,r}}.
\]
When $r$ owns $D$, the boundaries are the commits of $D$'s own release tags,
and the windows are those of \cref{def:window}. When $r$ does not own $D$,
the boundaries are the positions $D$'s releases prove in $r$, and only these
prove them: a control release checkpoint whose release-commit message names
$D$'s exact tag, changes $D$'s gitlink to the commit peeled from that tag, and
pins the revision of each other repository the release incorporated
(\spec{27.6}); an explicit \code{repositoryBaselines} tuple $(D,
\text{release tag}, r, \text{revision})$, with stable and prerelease tags as
separate entries (\spec{27.6}); or, in the linked peer topology, the route of
fleet links recorded hop by hop from the tree of the release commit that
names $D$'s tag (\spec{27.11}). A matching object identifier alone proves
nothing. Missing or ambiguous evidence is \diag{E333}, never a guess and
never a choice by date or proximity. Admission of a unit $u$ carried by a
commit of $r$ at a consumer $D$ remains $\cof{u} \in \exWinf{D}{r}$: $D$'s
last release had not yet incorporated the qualified revision. A source unit
resolves directly only packages of its own repository and propagates through
the one combined graph; delivery (\spec{13.4a}) compares $D$'s boundary and
the release commits of the provider in the provider's repository. A window
over $r$ is a function of $r$, $h_r$ and one boundary revision of $r$ and of
nothing else, so packages of any owners share an immutable train window over
$r$ exactly when their stable boundary there is the same qualified revision,
and a fresh window exactly when both boundaries are.
\end{definition}

The identity $\exWinf{D}{r} = \exafter{r}{s_{D,r}} \cap \exafter{r}{b_{D,r}}$
is the fresh-window identity of \cref{prop:marker-pass}, read per
repository; it is what lets \cref{sec:complexity} key every window by one
boundary revision. Newest-wins precedence between competing directives is
defined within one repository's DAG and within the control history; source
revisions of different repositories have no such order, and where the
existing rule requires a single winner among incomparable revisions the
engine fails with \diag{E334} rather than choosing by date, traversal order,
repository name or identifier (\spec{27.4}). An explicit control directive can
resolve the choice, evaluated against the gitlinks its own commit pinned, its
causal snapshot, so that it cannot rewrite work the control revision had not
yet observed; in the linked peer topology no repository addresses the whole
fleet and the conflict is refused outright.

\begin{definition}[recording order, \spec{27.7}, \spec{27.11}]
\label{def:exec-recording}
The combined publish sequence is that of \cref{def:order} across all
repositories, and publication remains non-atomic. Each successful package
receives its immutable release tag in its owning repository at the ordinary
tag-after-publish point of \cref{def:tagging}. Where a source release commit
and push are enabled, the source release is recorded durably first, and only
after its branch and tags are reachable from the configured remote may a
control gitlink or a peer's pin advance to it: a pushed pointer never names a
commit nobody else can fetch. A control checkpoint, when enabled and
non-empty, follows the source record and never precedes it; no empty commit
is forced and no checkpoint is required. A failure to record or durably push
a tag is \diag{E335}; prior successes remain successful and dependents are
protected by \cref{def:blocking}. If the source record succeeds and the
control checkpoint then fails, \diag{E335} names the repository, the full
revision and the tag and requires explicit repair; the next run remains
\diag{E330} while the checkout and the committed gitlink are unpinned, and
the engine neither auto-commits the repair, nor forces a checkpoint, nor
republishes the source, nor infers success from the worktree. A successful
source tag is never deleted, moved or duplicated to make a checkpoint look
atomic, and no rollback is inferred. In the linked peer topology the route
evidence is settled by ordinary commits before publication, deepest hop
first, all or nothing per package (\diag{E333} otherwise), and each
settlement is an admitted transition of the repository it was written in.
The fleet lock covers the control repository, every active source and every
standalone package as one fleet (\diag{E336} otherwise), is taken in one
stable total order over lock identities, and is released in reverse.
\end{definition}

\paragraph{Bounds.}
\spec{27.9} requires the engine to index the control gitlink transitions
once per fixed control snapshot, not once per consumer or provider pair; to
store each parsed (repository, commit) record once; to share window
membership only under the keys of \cref{def:exec-windows}; and to bound its
reachability and ancestry caches, because a cache of every queried pair can
grow quadratically in one repository's history even when that history is
walked once. With the notation of \cref{sec:complexity}, input traversal is
$O(G + \sum_q (H_q + A_q))$, window reachability is $O(\sum_q (H_q + A_q)\cdot
\mathrm{bw}(m_q) + I_w)$ with one marker pass per repository, and in the linked
peer topology boundary resolution is $O(Q + X \cdot Y)$ lookups over $B$ tree
reads, with settlement in at most $Q - 1$ commits and at most two in a star.
Publication revalidation adds $O(\sum_q J_q (1 + T_q))$ comparison work, with
$J_q$ the fleet or package checks repository $q$ takes part in and $T_q$ its
relevant tag refs, and retaining every package's input closure costs memory
linear in the total number of repository memberships across those closures,
at most $P \cdot Q$. The profile makes no globally linear end-to-end claim and permits
no synthetic ancestry shortcut.

\begin{remark}[what this subsection does not prove]
\label{rem:exec-multihistory}
\Cref{lem:window}, and every theorem of \cref{sec:guarantees} that rests on
it, is stated for one history with one \Head{}. Under \cref{def:exec-windows}
a consumer has one window per repository its plan reads, each measured from a
boundary its own release proves there, and a release of the consumer moves
every one of those boundaries at once only when the evidence of
\cref{def:exec-recording} is written. Restating the retry invariant over the
fleet snapshot and proving the eight guarantees per repository is a separate
piece of work; this report defines the objects, and \spec{27} promises that
the single-history rules apply unchanged inside each repository.
\end{remark}
